\section{Limitaciones y trabajo futuro}
\label{sec:limitaciones}

Un modelo se defiende mejor declarando sus límites que escondiéndolos. Esta
sección separa lo que se omitió a sabiendas, con su justificación, de lo que
se detectó durante la verificación o en la revisión externa y sigue
abierto, y termina con los datos del dron real que hacen falta para cerrar
varios de esos puntos.

\subsection{Omisiones del modelo físico}

\begin{itemize}
\item \textbf{Dinámica eléctrica del motor y reacción por aceleración del
rotor.} El propulsor de primer orden no satura la corriente, por lo que su
aceleración inicial ante un escalón es irreal y el término $-\dot{\vect h}_r$
se omite (\cref{sec:reaccion}). Un modelo con resistencia de devanado,
constante $K_t$ y límite de corriente permitiría incluirlo y daría además
una medida directa de calentamiento, en lugar del sustituto
\cod{motor\_hf\_pct}.
\item \textbf{Constante de tiempo del motor.} $T_m=20$~ms viene de
CopterSim, no de los motores del proyecto, y es la misma en todo punto de
operación, con carga o sin ella, al acelerar y al frenar, sin límite de
pendiente. El lazo de tasa es más sensible a $T_m$ que a los retardos
omitidos: con 30~ms el margen de fase baja de 58.7 a 48.1$^\circ$
(\cref{sec:integracion}).
\item \textbf{Batería.} Voltaje constante: $C_R$ no cae durante el vuelo y
el acelerador de hover no deriva. CopterSim trae un modelo de batería que
podría portarse.
\item \textbf{Aerodinámica del rotor.} No hay \emph{blade flapping}, arrastre
inducido ni pérdida de empuje en vuelo de avance, que a velocidades de
varios m/s producen fuerzas horizontales y momentos de cabeceo
proporcionales a la velocidad. Tampoco hay anillo de vórtices en descensos
rápidos. El arrastre vertical es un supuesto y el momento del arrastre por
la separación de 0.12~m entre centro aerodinámico y de masa que declara
CopterSim no se modela.
\item \textbf{Efecto suelo.} Escala el empuje de cada rotor pero no su par
de arrastre. Es un supuesto de modelado, no una consecuencia del modelo de
efecto suelo elegido.
\item \textbf{Forma del empuje.} $T=C_T\Omega^2$ puro, frente al ajuste
$b\Omega^2+a\Omega+c$ que Quitian Niño \cite{quitian22} encontró
significativo en banco.
\item \textbf{Inercia y carga.} Tensor diagonal; la carga útil no agrega
productos de inercia y su desplazamiento vertical del centro de masa no
produce brazo para el empuje.
\item \textbf{Turbulencia.} Filtros de primer orden en los tres ejes,
alineados con NED (\cref{sec:perturbaciones}).
\item \textbf{Suelo.} Restricción simple sin fricción ni rebote.
\end{itemize}

\subsection{Omisiones de la cadena digital}

\begin{itemize}
\item \textbf{ESC.} Por defecto responde de forma instantánea y continua a
$\sigma$: no hay frecuencia de PWM, protocolo (PWM, OneShot, DShot) ni
filtrado de entrada. Con PWM de 50~Hz el margen de fase del lazo de tasa
caería a 39.6$^\circ$ y la ganancia admisible a 2.4 veces
(\cref{tab:margenes}). Las opciones \cod{esc\_rate\_hz} y
\cod{esc\_pulse\_delay\_s} permiten estudiarlo, pero el valor real del dron
del proyecto no se conoce todavía.
\item \textbf{Retardo de cómputo.} Por defecto el mando se aplica en el
mismo periodo en que se calcula; \cod{compute\_delay\_samples} agrega uno o
más periodos. Su efecto sobre los márgenes está en la
\cref{sec:integracion}.
\item \textbf{Resolución del acelerador.} $\sigma$ es continuo por defecto.
Con el LEDC del ESP32 a 50~Hz y 10 bits quedan unos 51 escalones útiles,
que pueden generar ciclos límite; \cod{throttle\_levels} permite
cuantizarlo.
\item \textbf{Temporización ideal.} El lazo es perfectamente periódico: no
hay \emph{jitter} de FreeRTOS (interrupciones de la radio NRF24 o del WiFi,
lecturas I2C), ni desfase entre la tasa de salida de la IMU y el lazo de
500~Hz, ni el filtro digital interno de la IMU. El retardo de 2~ms de la IMU
es un supuesto sin fuente citada.
\item \textbf{Barómetro.} Un ruido de 0.17~m corresponde a un modo de alta
resolución del BMP180, cuya conversión limita la tasa a unos 13--33~Hz; la
simulación usa 40~Hz. Falta contrastarlo con la hoja de datos y con la
configuración del firmware.
\item \textbf{Enlace de posición.} Además del retardo de 10~ms del sensor,
\cod{pos\_link\_delay\_s} suma la latencia del enlace (captura de movimiento
o radio, típicamente 5--30~ms). Vale 0 por defecto porque depende del
sistema de posicionamiento que se use con el dron real.
\end{itemize}

Los valores por defecto de las opciones nuevas (\cod{compute\_delay\_samples},
\cod{esc\_rate\_hz}, \cod{esc\_pulse\_delay\_s}, \cod{throttle\_levels},
\cod{pos\_link\_delay\_s}) conservan el
comportamiento con el que se identificaron los rezagos y se afinaron los
controladores, para que las cifras de este documento sigan siendo
reproducibles. Una sintonía hecha con ellas desactivadas es optimista en
margen de fase. Antes de llevarla al dron conviene repetir el escalón de
tasa con los valores reales del ESC y del firmware.

\subsection{Hallazgos de la verificación}

Varios hallazgos de la verificación ya se corrigieron en el núcleo y este
documento describe el estado corregido: el pasa-bajos del acelerómetro antes
del AHRS, que elimina el sesgo de inclinación por rectificación de la
vibración (de \val{rectmadC} a \val{rectmadF}$^\circ$ con Madgwick); el
arranque en el aire coherente (líneas de retardo llenas, AHRS alineado,
covarianza de navegación en régimen); el filtro que se rediseña en vuelo
sin transitorio; los valores por defecto de $k_d$, del filtro del
giroscopio, del EKF y de $\beta$. La revisión externa de octubre de 2026
llevó a corregir además: la compensación de aceleración, que ya no se aplica
sin correcciones horizontales de posición absoluta, entra en rampa, se
recorta y pasa por el mismo filtro que el acelerómetro
(\cref{sec:leishman}); la corrección magnética de Madgwick y Mahony, que
ahora solo actúa sobre el rumbo; la base de tiempo con razones no enteras
entre física y control (\cref{sec:integracion}); la línea de retardo de los
sensores, que ahora admite retardo nulo y fraccional y arranca cebada; la
validación de nombres de estimadores y de ganancias, que antes ignoraba sin
aviso una clave mal escrita; el tope del integrador de Mahony; el registro,
que ya no descarta columnas que aparecen a mitad de corrida; y una
regularización adaptativa en el iLQR del NMPC. Quedan abiertos:

\begin{itemize}
\item \textbf{Compensación de aceleración.} Con sensor de posición sigue
siendo autorreferente entre correcciones y queda un desfase de unos 17~ms
entre las dos señales que se restan (\cref{sec:leishman}). Sin sensor de
posición queda desactivada, así que en ese modo el AHRS vuelve a
equivocarse durante las aceleraciones sostenidas. La solución de fondo es
un filtro de navegación con el error de actitud como estado.
\item \textbf{Error de actitud en ángulos de Euler.} El lazo de ángulo resta
ángulos de Euler. Es válido con el límite de 30$^\circ$, pero se degrada en
maniobras agresivas y es singular cerca de $\pm90^\circ$ de cabeceo; PX4
usa el error de cuaternión.
\item \textbf{Madgwick sin estimación de sesgo.} El término de compensación
de deriva del artículo original no está implementado; la comparación con
Mahony y el EKF, que sí estiman el sesgo, debe leerse con eso en mente.
\item \textbf{Retardo de la cadena en la sintonía.} Las ganancias por
defecto se calculan con la cadena nominal ($\theta=16.0$~ms); la cadena
configurada, con filtros, cómputo y ESC, solo se usa con
\cod{gains\_from\_chain} (\cref{sec:control}). Ninguna de las dos incluye
el filtro del derivativo ni la fase exacta del motor.
\item \textbf{Barridos de estimadores anteriores.} Las tablas de
\cod{docs/BARRIDOS.md} se obtuvieron con la cadena de sombra anterior, que
no era la que se vuela (\cref{sec:barridos}); habría que repetirlos.
\item \textbf{Rezagos y políticas sin regenerar.} \cod{DEFAULT\_LAGS}, los
pesos del MPC y del NMPC y las políticas de RL se obtuvieron con el núcleo
anterior (Madgwick, sin filtro de acelerómetro); no se volvieron a
identificar, afinar ni entrenar (\cref{sec:posicion}). Re-identificar el
sustituto hoy da otros rezagos, así que la receta de \cod{docs/RL.md} no
reproduce exactamente los archivos de \cod{models/}.
\item \textbf{Ciclo límite del PID con EMI en escalones.} Incluso con
$K_{i,xy}=2$ (\cref{sec:posicion}).
\item \textbf{Exponente de la EMI.} $\sigma^{1.5}$ no se deriva del modelo de
propulsor (\cref{sec:perturbaciones}).
\item \textbf{Fusión del barómetro con el sensor de posición activo.} Inocua
en la práctica, pero no intencional (\cref{sec:estimadores}).
\end{itemize}

\subsection{Alcance del banco de controladores}

\begin{itemize}
\item \textbf{Tres semillas.} Bastan para ver tendencias consistentes, no
para probar diferencias de milímetros (\cref{sec:lectura_banco}).
\item \textbf{Afinación sobre el conjunto de prueba.} PID, MPC y NMPC se
afinaron con la semilla 1 de la figura en 8 y el círculo, que forma parte
del banco, y el integral del PID se eligió mirando el caso con EMI. El RL no
tuvo esa ventaja.
\item \textbf{Aprendizaje por refuerzo.} Un solo entrenamiento PPO por
política, cortado cuando la curva todavía bajaba; previsión, observación,
salida y límite de inclinación distintos de los del MPC; validación con
trayectorias que pertenecen a la familia de entrenamiento; y un residuo
elegido por costo y no por error. Las conclusiones valen para estas
políticas, no para el RL en general.
\end{itemize}

\subsection{Datos pendientes del dron real}

Cuatro mediciones cerrarían buena parte de lo anterior:
\begin{enumerate}
\item protocolo y frecuencia de actualización del ESC del dron (PWM,
OneShot o DShot, y su tasa);
\item constante de tiempo del motor medida en banco, con escalones de
acelerador y tacómetro óptico, en subida y en bajada;
\item modo y tasa reales del BMP180 en el firmware;
\item sistema de posicionamiento disponible (captura de movimiento, RTK o
ninguno) y su latencia, que decide si la compensación de aceleración puede
usarse y cuánto retardo tiene la posición.
\end{enumerate}

\subsection{Trabajo futuro}

El paso que más valor agrega es medir el dron real del proyecto y
reemplazar los parámetros de la literatura: masa y brazo, $C_T$ y $C_M$ en
banco, la recta del ESC y $T_m$ con escalones, inercias con péndulo
compuesto y trifilar (el protocolo de Quitian Niño \cite{quitian22}, de esta
misma universidad, ya describe esos bancos), y el espectro de vibración en
hover. Con esos números y los datos pendientes de arriba, el flujo que este
documento verifica (identificar los rezagos, derivar ganancias SIMC, afinar
en simulación con barridos y validar con el banco) se repite sobre el
vehículo real. Para los controladores de aprendizaje el orden es el
contrario del que sugería la edición anterior: primero mejorar la política
(observaciones normalizadas, la misma previsión que el MPC, salida como
corrección sobre la aceleración de referencia, selección por error de
validación y varias semillas de entrenamiento) y después, con una ablación
de un factor a la vez, decidir qué le falta al sustituto.
